\fichatitulo{51 · META-AVALIAÇÃO}{COLING · 2025}{GroUSE}
{\small Muller et al.. \textit{GroUSE: A Benchmark to Evaluate Evaluators in Grounded Question Answering}. COLING · 2025. \href{https://aclanthology.org/2025.coling-main.304/}{Fonte consultada}.\par}
\campo{Problema e objetivo}
Como testar se um avaliador de respostas fundamentadas mede realmente a qualidade que pretende medir?
\campo{Principal contribuição · síntese interpretativa}
Propõe um benchmark de 144 testes unitários para avaliar avaliadores automáticos de grounded QA, deslocando o foco da correlação agregada para falhas diagnósticas.
\campo{Método / natureza da evidência}
Conjunto controlado de respostas com alterações de suporte, contradição e relevância; comparação de avaliadores baseados em LLM e métricas automáticas.
\campo{Resultado ou argumento principal}
Correlação com um juiz forte não garante cobertura de casos-limite; avaliadores podem falhar sistematicamente em distinções que o protocolo exige.
\campo{Excertos-chave · seleção literal}
\begin{itemize}
\excerto{1}{a benchmark to evaluate evaluators in grounded question answering}{Título.}
\end{itemize}
\campo{Limitações e análise crítica}
Os testes unitários não substituem validação humana no domínio parlamentar. A transferência depende das dimensões efetivamente representadas na bateria de perguntas.
\campo{Aproveitamento na dissertação · proposta}
Justificar testes de calibração e casos-limite para o LLM-as-a-judge, além da correlação global.
\registro{Fonte consultada nesta preparação: PDF publicado; consulta focal do resumo, método e conclusões. Chave: \texttt{mullerEtAl2025grouse}.}
\clearpage
